\newpage
\section{Results and Discussion}
\label{Section:result}



\subsection{Pyrolysis for Biochar Production}
\setresponsible{Regen Wijaya}
\label{subsec:results-pyrolysis}

Table~\ref{tab:yield_summary} summarizes the data obtained from the pyrolysis and biochar yield calculated using Equation~\ref{eq:biochar yield}. Since the amount of biochar is quite small, the standard deviation was also calculated for each yield. 


\begin{table}[htbp]
\centering
\caption{Yields in biochar production from French and Algerian digestate via pyrolysis at different temperatures}
\label{tab:yield_summary}
\begin{tabular}{llccc}
\toprule
\multirow{2}{*}{\thd{Sample}} & \multirow{2}{*}{\thd{$T_\mathrm{pyr}$ (\textdegree C)}} & \multicolumn{2}{c}{\thd{Mass (g)}} & \multirow{2}{*}{\thd{Yield (\%)}} \\
\cmidrule(lr){3-4}
 & & \thd{Digestate} & \thd{Biochar} & \\
\midrule
\multirow{3}{*}{French digestate}   & 400 & $14.068 \pm 0.041$ & $6.436 \pm 0.023$ & $45.75 \pm 0.05$ \\
                                    & 500 & $14.044 \pm 0.005$ & $6.056 \pm 0.006$ & $43.12 \pm 0.05$ \\
                                    & 600 & $15.122 \pm 0.434$ & $6.379 \pm 0.193$ & $42.18 \pm 0.07$ \\
\midrule
\multirow{3}{*}{Algerian digestate} & 400 & $2.402 \pm 0.008$  & $0.899 \pm 0.006$ & $37.42 \pm 0.32$ \\
                                    & 500 & $2.406 \pm 0.006$  & $0.827 \pm 0.003$ & $34.37 \pm 0.14$ \\
                                    & 600 & $2.405 \pm 0.007$  & $0.805 \pm 0.016$ & $33.45 \pm 0.58$ \\
\bottomrule
\end{tabular}

\tabnote{``$T_\mathrm{pyr}$'' refers to the pyrolysis temperature at which the
biochars were produced.}
\end{table}

Across the pyrolysis temperature range, the French digestate exhibited a 8--9\% higher biochar yield, decreasing from 45.7\,\% at 400\,$^{\circ}$C to 42.2\,\% at 600\,$^{\circ}$C, whereas the Algerian digestate yielded between 37.4\,\% and 33.5\,\% over the same temperatur range, as shown in Table~\ref{tab:yield_summary}. This difference may be attributed to different feedstock compositions between both digestates rather than the pyrolysis process because both digestates were derived from slaughterhouse waste and pyrolyzed under identical conditions. The French digestate may have a larger fraction of material retained in the solid phase, which points to differences in the organic and mineral balance of the two digestates arising from their respective origins and upstream digestion \cite{kujawska-2026,sarkar-2024}. This interpretation was addressed further in Section~\ref{sec:characterisation} based on the elemental and ash data, where the mineral content retained in the biochar is directly quantified.

\vspace{0.3em}

Figure~\ref{fig:pyrolysis-yield} also provides a visual overview of how the biochar yield changed across the investigated temperature range for both digestates. As shown in Figure~\ref{fig:pyrolysis-yield}, the biochar yields from both digestates followed the same trend, with a lower standard deviation for the entire French digestate below 0.1 than for the Algerian digestate in Table~\ref{tab:pyrolysis-yield-full}. Both biochar yields decreased steadily as the pyrolysis temperature increased, amounting to about 3.6 percent for the French biochar and 4.0 for the Algerian biochar between 400 and 600\,$^{\circ}$C.

\begin{figure}[ht]
\centering
\begin{tikzpicture}
\begin{axis}[
    thesisline,
    title={Biochar production at different pyrolysis temperatures},
    xlabel={Pyrolysis temperature ($^{\circ}$C)},
    ylabel={Biochar yield (\%)},
    xtick={400,500,600},
    xmin=370, xmax=630,
    ymin=30, ymax=48,
]
\addplot[s primary,
    error bars/.cd, y dir=both, y explicit]
  coordinates {
    (400,45.75) +- (0,0.05)
    (500,43.12) +- (0,0.05)
    (600,42.18) +- (0,0.07)
  };
\addlegendentry{French digestate}
\addplot[s secondary,
    error bars/.cd, y dir=both, y explicit]
  coordinates {
    (400,37.42) +- (0,0.32)
    (500,34.37) +- (0,0.14)
    (600,33.45) +- (0,0.58)
  };
\addlegendentry{Algerian digestate}
\end{axis}
\end{tikzpicture}
\caption{Biochar yield from pyrolysis at different temperatures for the French and
Algerian digestate.}
\label{fig:pyrolysis-yield}
\end{figure}


\newpage

The downward trend points to the expected release of volatile matter at higher temperatures since the thermally less stable fractions are removed and a more carbonized residue remains \cite{zhang-2014,zhang-2022,224pyrolysistemp}. The decline may also suggest that the high mineral fraction in digestate-derived feedstocks is largely retained in the solid during pyrolysis, making the yield responds less steeply to temperature than is typical for lignocellulosic biomass (Section~\ref{subsubsec:digestate-pyrolysis relevance}) \parencite{kujawska-2026, wang-2023}. This supports the assumption that the differences between the starting materials and the temperatures are due to actual compositional effects rather than measurement variations.

\vspace{0.3em}

The biochar yields of this work reproduced similar values with little differences for Algerian biochars compared to biochar yield of the preceding project work, as shown in Table~\ref{tab:projektarbeit-yield}, as they were produced with the same conditions and procedures. The Algerian digestate, which was pyrolyzed at 400 °C and 500 °C in both studies, yielded identical amounts of biochar. However, the Algerian digester residue that was pyrolyzed at 600 °C in the previous project yielded 2\% less biochar compared to the Algerian digestate pyrolyzed at 600 °C in the present study, which could be attributed to human error rather than to the composition of the digester residue. 

\begin{table}[ht]
\centering
\caption{Biochar yield of the three digestates examined in the
preceding project work \parencite{MartinezHuete2024Projektarbeit}.}
\label{tab:projektarbeit-yield}
\begin{tabularx}{\textwidth}{@{}CCCC@{}}
\toprule
\thd{Pyrolysis temperature (°C)} & \thd{Slaughterhouse yield (\%)} \\
\midrule
400 & 37 \\
500 & 34 \\
600 & 31 \\
\bottomrule
\end{tabularx}
\end{table}



\newpage

\subsection{\texorpdfstring{H$_2$O}{H2O}-Activation}
\setresponsible{Surya Salim}
\label{subsec:results-activation}

Figure~\ref{fig:activation-massloss} shows the activation mass loss measured for the two water doses (0.5\,mL\,g$^{-1}$ and 0.1\,mL\,g$^{-1}$). All four data series show an overall decrease with increasing the pyrolysis temperature. On the moisture-corrected basis, the mass loss decreased, from 17.9\,\% at 400\,$^{\circ}$C to 11.1\,\% at 600\,$^{\circ}$C, and the 2\,g scale-up runs followed the same descending order (20.6\,\%, 15.0\,\% and 9.8\,\%; Table~\ref{tab:activation-full}). The uncorrected wet basis further supports this downward trend, as shown by the dashed line in Figure~\ref{fig:activation-massloss}.

\begin{figure}[ht]
\centering
\begin{tikzpicture}
\begin{axis}[
    thesisline,
    title={Activation mass loss of activated French biochar},
    xlabel={Pyrolysis temperature ($^{\circ}$C)},
    ylabel={Activation mass loss (\%)},
    xtick={400,500,600},
    xmin=370, xmax=630,
    ymin=0, ymax=24,
]
% Hue = water dose (blue 0.5 mL, vermillion 0.1 mL);
% solid/filled = moisture-corrected, dashed/open = wet basis.
\addplot[s primary]
  coordinates {(400,17.92) (500,15.86) (600,11.05)};
\addlegendentry{0.5 mL, corrected}
\addplot[s secondary]
  coordinates {(400,20.76) (500,21.65) (600,10.29)};
\addlegendentry{0.1 mL, corrected}
\addplot[s primary, s open]
  coordinates {(400,8.91) (500,7.92) (600,3.96)};
\addlegendentry{0.5 mL, wet}
\addplot[s secondary, s open]
  coordinates {(400,15.84) (500,7.92) (600,4.00)};
\addlegendentry{0.1 mL, wet}
\end{axis}
\end{tikzpicture}
\caption{Mass loss of the French biochar during activation at both water dosages (0.1 and
0.5\,mL\,g$^{-1}$) as a function of pyrolysis temperature.}
\label{fig:activation-massloss}
\end{figure}

This trend proposes that biochar produced at higher pyrolysis temperatures yields a more carbonized and thermally stable solid in which the volatile fractions have already been released (Section~\ref{subsubsec:Effect of Temperature on Biochar}), so less material remains available for removal during activation. The pattern is actually the counterpoint to pyrolysis line graph in Section~\ref{subsec:results-pyrolysis}. Volatile matter removed during pyrolysis is no longer present during activation, thereby, the more extensively devolatilized high-temperature biochars are, the less additional loss there is under activation. Residual devolatilization may also contribute substantially to the mass loss during activation. All three biochars were produced below the 700\,$^{\circ}$C activation set point, so reheating releases volatiles that have not yet been removed through pyrolysis beforehand. The effect is largest for 400\,$^{\circ}$C and smallest for the 600\,$^{\circ}$C biochar.

\vspace{0.3em}

As mentioned in section \ref{subsubsec:physical-activation}, steam gasification is the pore-forming reaction, which depends heavily on a steam supply, activation temperatures and activation time. Mass loss on its own cannot separate steam gasification from residual devolatilization, so it cannot settle whether the 0.5\,mL\,g$^{-1}$ charge reached the greater carbon conversion despite its smaller total loss. That question is left to the BET results. The 0.5\,mL\,g$^{-1}$ dose was retained as the standard for its more reliable wetting of the char and its lower mass loss.

\begin{table}[ht]
\centering
\caption{Activation mass loss from CO$_2$ activation for the three digestate-derived biochars of
the preceding project work
\parencite{MartinezHuete2024Projektarbeit}.}
\label{tab:co2-activation-projektarbeit}
\begin{tabular}{lccc}
\toprule
\thd{Feedstock} & \thd{Pyrolysis temp.\ ($^{\circ}$C)} & \thd{Residual char (g)}
& \thd{Mass loss (\%)} \\
\midrule
\multirow{3}{*}{Slaughterhouse digestate}& 400 & 0.85 & 15 \\
                                         & 500 & 0.93 & 7  \\
                                         & 600 & 0.98 & 2  \\
\bottomrule
\end{tabular}
\end{table}


\begin{figure}[ht]
\centering
\begin{tikzpicture}
\begin{axis}[
    thesisline,
    title={Steam vs $CO_2$ activation mass loss},
    xlabel={Pyrolysis temperature of biochar ($^{\circ}$C)},
    ylabel={Activation mass loss (\%)},
    xtick={400,500,600},
    xmin=370, xmax=630,
    ymin=0, ymax=24,
]
\addplot[s primary]
  coordinates {(400,17.92) (500,15.86) (600,11.05)};
\addlegendentry{Steam, 0.5\,mL\,g$^{-1}$ (corrected)}
\addplot[s secondary]
  coordinates {(400,20.76) (500,21.65) (600,10.29)};
\addlegendentry{Steam, 0.1\,mL\,g$^{-1}$ (corrected)}
\addplot[s reference]
  coordinates {(400,15) (500,7) (600,2)};
\addlegendentry{CO$_2$, slaughterhouse digestate (Algerian)}
\end{axis}
\end{tikzpicture}
\caption{Activation mass loss versus pyrolysis temperature for the French
biochar under steam activation and for the slaughterhouse-digestate biochar under CO$_2$ activation
\parencite{MartinezHuete2024Projektarbeit}.}
\label{fig:co2-vs-steam}
\end{figure}

The CO$_2$-activation of the preceding project work offers a direct point of comparison since it also used the same reactor, temperature, and residence time as the steam-activation runs here (700\,$^{\circ}$C, 1\,h). These two slaughterhouse-digestates, though they come from different sources, could provide insight into this discussion. Since no water is introduced during CO$_2$ activation, the reported before/after loss is assumed to be under low-moisture conditions and can be placed alongside the corrected steam values without adjustment. 

\vspace{0.3em}

Based on this, the steam-activated char loses more mass than the CO$_2$-activated char at every pyrolysis temperature (Figure~\ref{fig:co2-vs-steam}; Table~\ref{tab:co2-activation-projektarbeit}). The widest gap could be seen at 600\,$^{\circ}$C, where steam removes about 10--11\,\% compared to roughly 2\,\% for CO$_2$. Steam tends toward a more extensive distribution of micropores and also mesopores due to pore widening. It also adds more oxygen-containing surface groups, whereas CO$_2$ works mainly within the micropore range, as mentioned in \ref{subsubsec:physical-activation}. 

\subsection{Sulfur loading}
\setresponsible{Regen Wijaya}
\label{subsec:results-sulfur}
The mass gain of each biochar after the sulfur loading process was calculated from the difference between the mass before and after the sulfur loading. The sulfur adsorption capacity of each biochar was then calculated using Equation~\ref{Equation sulfur loading} and summarized together with the collected data in Table~\ref{tab:sulfur_loading}. 

\begin{table}[H]
\centering
\caption{Sulfur adsorption capacity of the steam-activated French biochars.}
\label{tab:sulfur_loading}


\tabwide
\begin{tabular}{lcccccc}
\toprule

\thd{Sample}
& \thd{Water}
& \thd{$T_\mathrm{pyr}$}
& \multicolumn{2}{c}{\thd{Mass of biochar (g)}}
& \thd{Mass gain}
& \thd{Sulfur adsorption capacity}
\\

& \thd{(mL)}
& \thd{($^\circ$C)}
& \thd{Before}
& \thd{After}
& \thd{(g)}
& \thd{(~mg\,g$^{-1}$)}
\\

\midrule

\multirow{6}{*}{French Biochar}
& \multirow{3}{*}{0.1}
& 400
& 0.829
& 0.875
& 0.045
& 54.504
\\

&
& 500
& 0.829
& 0.882
& 0.053
& 64.302
\\

&
& 600
& 0.904
& 0.946
& 0.041
& 45.671
\\

\cmidrule(lr){2-7}

&
\multirow{3}{*}{0.5}
& 400
& 0.900
& 0.936
& 0.037
& 40.912
\\

&
& 500
& 0.932
& 0.966
& 0.034
& 36.373
\\

&
& 600
& 0.923
& 0.963
& 0.040
& 43.883
\\

\bottomrule
\end{tabular}

\tabnote{``Water'' refers to the amount of water added during activation.}

\end{table}

Table~\ref{tab:sulfur_loading} shows that all biochars exhibited an increase in mass after the sulfur loading process. This indicates that elemental sulfur was adsorbed during the process and retained on the surface of the activated biochars. The biochars activated with 0.1 ml water dose presented a greater mass increase in the range of 0.04 to 0.05 g and an sulfur adsorption capacity in the range of 45 to 64 ~mg\,g$^{-1}$ compared to the biochars activated with a water dose of 0.5 ml, for which the mass increase ranged from 0.034 to 0.04 g and the adsorption capacities ranged from 36 to 44 ~mg\,g$^{-1}$. This difference is too small to carry a mechanism. Each dose was run once per pyrolysis temperature, and the spred of about 20\,mg\,g$^{-1}$ between the two doses is comparable to the disagreement between the gravimetric and the CHNS sulfur figures for the same samples, which reaches 16\,mg\,g$^{-1}$ at 400\,°C (Section~\ref{subsubsec:chns}). The activation data argue against a dose effect as well: on the moisture-corrected basis the 0.1\,mL dose removed more mass than the 0.5\,mL dose at 400 and 500\,°C (Table~\ref{tab:activation-full}), so the higher dose was not the harsher condition, and BET was run only on the 0.5\,mL samples, so there is no surface-area comparison between the doses.
%This can imply that steam activation with greater activation water quantity may lead to excessive gasification, in which micropores are converted into mesopores or macropores. As a result, the specific surface area to retain more elemental sulfur can be reduced, thereby decreasing the area available for adsorption. 

\vspace{0.3em}

Among the biochars activated with 0.1 ml water dose, the highest mass increase of approximately 0.053 g and the highest sulfur adsorption capacity of approximately 64.3 ~mg\,g$^{-1}$ were observed in the biochar pyrolyzed at 500 °C, followed by the biochar pyrolyzed at 400 °C and 600 °C. This may possibly be attributed to further possible carbonization at 600 °C, which can lead to a reduction in oxygen-containing functional groups \cite{Sajjadi2019} \cite{tomczyk-2020}, although higher pyrolysis temperatures are generally considered advantageous for higher sulfur adsorption capacity, as they promote pore formation and result in a larger specific surface area. 

%\begin{figure}[htbp]
%\centering
%\begin{tikzpicture}
%\begin{axis}[
%    width=1\textwidth,
%    height=0.62\textwidth,
%    xtick={400,500,600},
%    xticklabels={{400\,$^{\circ}$C},{500\,$^{\circ}$C},{600\,$^{\circ}$C}},
%    title={Adsorption capacity of French biochar},
%    ylabel={Adsorption capacity in ~mg\,g$^{-1}$},
%    xlabel={Pyrolysis temperature in $^\circ$C},
%    ymin=25, ymax=70,
%    xmajorgrids=true,
%    ymajorgrids=true,
%    grid style={draw=gray!25},
%    tick label style={font=\small},
%    label style={font=\small},
%    legend style={at={(0.5,-0.18)}, anchor=north, legend columns=2,
%                   font=\small, draw=gray!40,
%                  /tikz/every even column/.append style={column sep=0.4cm}},
%]
%\addplot[s primary]
% coordinates {(400,54.504) (500,64.302) (600,45.671)};
%\addlegendentry{Adsorption capacity, 0.1\,mL}
%\addplot[s secondary]
%  coordinates {(400,40.912) (500,36.373) (600,43.883)};
%\addlegendentry{Adsorption capacity, 0.5\,mL}
%\end{axis}
%\end{tikzpicture}
%\caption{Adsorption capacity of French biochar after sulfur loading, at two water activation doses (0.1 and 0.5\,mL) across pyrolysis temperatures.}
%\label{fig:sulfur-loading-capacity}
%\end{figure}

%\vspace{0.3em}

%\newpage

%Figure~\ref{fig:sulfur-loading-capacity} illustrates the adsorption capacity of French biochar samples after sulfur loading as a function of pyrolysis temperature. The samples activated using 0.1 and 0.5 mL water exhibited different behavior at different pyrolysis temperature. For the samples activated with 0.1 mL water, the adsorption capacity increases from the biochar produced at 400°C to  



%Figure~\ref{fig:sulfur_loading_capacity} visualizes the influence of pyrolysis temperature and activation-water dose on the sulfur-loading capacity. The two activation conditions show different temperature-dependent trends. For the samples activated with 0.1~mL water, the capacity increases from 400 to 500,$^{\circ}$C and then decreases considerably at 600,$^{\circ}$C, resulting in a clear maximum at 500,$^{\circ}$C. In contrast, the 0.5~mL series shows a slight decrease between 400 and 500,$^{\circ}$C followed by an increase at 600,$^{\circ}$C. Therefore, no simple monotonic relationship between pyrolysis temperature and sulfur-loading capacity can be observed.
\newpage
In contrast, among the biochars activated with a 0.5 ml water dose, the biochar pyrolyzed at 600 °C showed the highest increase of approximately 0.04 g and an sulfur adsorption capacity of 43.884 ~mg\,g$^{-1}$, while the lowest values were observed for the sample produced by pyrolysis at 500 °C. This implies that the effect of the pyrolysis temperature on the sulfur adsorption capacity may also be governed by the amount of water used during activation, as it influences the pore development and accessible internal surface for adsorption. 

%Although higher pyrolysis temperatures generally promote pore development and may increase the specific surface area of biochar, the highest average apparent adsorption capacity in this study was obtained for the biochar produced at 500 °C. This may indicate that the biochar produced at 500 °C retained a more favorable balance between accessible pore structure and reactive surface functional groups. At 600 °C, further carbonization may have reduced the availability of oxygen-containing functional groups involved in H₂S retention. However, because the adsorption performance also varied with the amount of activation water, the observed differences cannot be attributed to pyrolysis temperature alone.

%Pyrolysis temperature also significantly influences pore development and surface area. Elevated temperatures promote devolatilization and the formation of more developed porous structures, which are beneficial for adsorption-related applications such as wastewater treatment, gas purifi-cation, and contaminant removal. Higher-temperature biochars generally possess larger surface areas and improved adsorption performance due to increased aromaticity and pore formation 

\vspace{0.3em}

%Among the samples activated with 0.1 mL of water, the greatest mass increase was observed for biochar produced by pyrolysis at 500 °C. In contrast, among the samples activated with 0.5 mL of water, the greatest mass increase was observed for biochar produced by pyrolysis at 600 °C. 

%Overall, biochar activated with the lower water dose (0.1 mL) exhibited higher sulfur adsorption capacities (45.671–64.302 ~mg\,g$^{-1}$) compared to those activated with the higher dose (0.5 mL), which yielded capacities between 36.373 and 43.883 ~mg\,g$^{-1}$. This suggests that a lower water content during the activation process is more favorable for sulfur uptake, possibly because excess water occupies or obstructs the pore structure of the biochar, thereby reducing the surface area accessible for sulfur adsorption.

\begin{table}[H]
\centering
\caption{Sulfur adsorption capacity of the  CO$_2$-activated Algerian biochars.}
\label{tab:sulfur_loading_co2}


\begin{tabular}{lccccc}
\toprule

\thd{Sample}
&
\thd{($^\circ$C)}
&
\multicolumn{2}{c}{\thd{Mass of biochar (g)}}
&
\thd{Mass gain}
&
\thd{Sulfur adsorption capacity}
\\

&
\thd{($^\circ$C)}
&
\thd{Before}
&
\thd{After}
&
\thd{(g)}
&
\thd{(~mg\,g$^{-1}$)}
\\

\midrule

\multirow{3}{*}{%
\begin{tabular}[c]{@{}l@{}}
Slaughterhouse\\
digestate\\
(Algeria)
\end{tabular}}
&
400
&
0.843
&
0.864
&
0.021
&
24.428
\\

&
500
&
0.950
&
0.965
&
0.016
&
16.426
\\

&
600
&
0.976
&
0.988
&
0.012
&
12.298
\\

\bottomrule
\end{tabular}

\end{table}

%\begin{figure}[htbp]
%\centering
%\begin{tikzpicture}
%\begin{axis}[
%    width=0.9\textwidth,
%    height=0.55\textwidth,
%    xbar,
%    bar width=16pt,
%    symbolic y coords={400,500,600},
%    ytick=data,
%    yticklabels={{400\,$^{\circ}$C},{500\,$^{\circ}$C},{600\,$^{\circ}$C}},
%    y dir=reverse,
%    title={Adsorption capacity of Algerian biochar},
%    xlabel={Adsorption capacity in ~mg\,g$^{-1}$},
%    ylabel={Pyrolysis temperature in $^\circ$C},
%     xmin=0, xmax=30,
%    xtick distance=10,
%    xmajorgrids=true,
%    grid style={draw=gray!25},
%    tick label style={font=\small},
%    label style={font=\small},
%    nodes near coords={\pgfmathprintnumber[fixed,precision=1]{\pgfplotspointmeta}},
%    every node near coord/.append style={font=\scriptsize, text=black, anchor=west},
%]
%\addplot+[b primary]
%  coordinates {(24.428,400) (16.426,500) (12.298,600)};
%\end{axis}
%\end{tikzpicture}
%\caption{Adsorption capacity of Algerian biochar activated using CO$_2$ from previous work.}
%\label{fig:adsoprtion capacity algeria biochar preivous}
%\end{figure}

Table~\ref{tab:tab:sulfur_loading_co2}) shows that the sulfur adsorption capacity of Algerian biochars activated with CO$_2$ decreased as the pyrolysis temperature increased, with the highest sulfur adsorption capacity of approximately 24 ~mg\,g$^{-1}$ observed for biochar pyrolyzed at 400 °C, suggesting that CO$_2$ activation is favored in biochar pyrolyzed at lower temperatures. 

\vspace{0.3em}

Compared to the sulfur adsorption capacity of CO$_2$-activated biochars from preceding work, the adsorption capacities of biochars in this work are 2-3 times higher. Although the CO$_2$ activation is theoretically associated with the development of narrow micropores, steam activation not only promotes the formation of micropores but also allows for the formation of mesopores and macropores through pore widening, which may facilitate the transport pathway of gas mixture of hydrogen sulfide and oxygen to the active adsorption sites \cite{Sajjadi2019}. The fewer presence of oxygen-containing functional groups such as phenolic and carbonyl groups during CO$_2$-activation may also contribute to the lower sulfur adsorption capacity of the biochar, since they provide reactive sites for the adsorption \cite{Sajjadi2019}. This suggests that the sulfur adsorption capacity may be influenced not only by the surface area but also by the pore size distribution and the oxygen-containing functional groups.

\newpage
\subsection{Characterization of biochar}
\label{sec:characterisation}

\subsubsection{Calorific value}
\setresponsible{Surya Salim}
\label{subsubsec:calorific}

\begin{figure}[ht]
\centering
\begin{tikzpicture}
\begin{axis}[
    name=maxis,
    thesisbar,
    bar width=13pt,
    enlarge x limits=0.18,
    symbolic x coords={Digestate,400,500,600},
    xticklabels={Digestate,{400\,$^{\circ}$C},{500\,$^{\circ}$C},{600\,$^{\circ}$C}},
    ylabel={Gross calorific value (MJ\,kg$^{-1}$)},
    ymax=32,
    legend columns=3,
    legend style={at={(0.5,-0.16)}, anchor=north, legend columns=3,
                  font=\small, draw=tcGridGrey, fill=white,
                  /tikz/every even column/.append style={column sep=0.5cm}},
]
\addplot+[b primary,
          error bars/.cd, y dir=both, y explicit]
  coordinates {(Digestate,16.41)+-(0,0.13) (400,14.16)+-(0,0.43)
               (500,14.17)+-(0,0.25) (600,14.11)+-(0,0.06)};
\addlegendentry{HHV, ash measured}
\addplot+[b emphasis]
  coordinates {(Digestate,21.20) (400,25.07) (500,27.88) (600,28.23)};
\addlegendentry{HHV, ash-free}
\addlegendimage{line legend, s reference}
\addlegendentry{Ash content}
\end{axis}
\begin{axis}[
    at={(maxis.south east)}, anchor=south east,
    thesisaxis right,
    symbolic x coords={Digestate,400,500,600},
    ylabel={Ash content (\%)},
    ymin=0, ymax=60,
]
\addplot[s reference,
         error bars/.cd, y dir=both, y explicit]
  coordinates {(Digestate,22.6)+-(0,0.6) (400,43.5)+-(0,2.0)
               (500,49.2)+-(0,0.9) (600,50.0)+-(0,0.6)};
\end{axis}
\end{tikzpicture}
\caption{Gross calorific value and ash content of the French digestate and its
biochars.}
\label{fig:calorific-value}
\end{figure}

The gross calorific value of the slaughterhouse digestate from France was 16.41\,MJ\,kg$^{-1}$. After pyrolysis, these values are lower for all three biochars at about 14.1--14.2\,MJ\,kg$^{-1}$ (Figure~\ref{fig:calorific-value}; Table~\ref{tab:calorific-full}). The calorific values of the biochars are visibly lower than that of the raw digestate itself. There are also minimal differences between pyrolysis temperatures, which might suggest that temperature has little effect on calorific value — but the ash trend argues otherwise. Ash was determined gravimetrically from the residue of the same combustion runs (Equation~\ref{eq:ash-content}, \S~\ref{subsubsec:methods-ash}). Bomb calorimetry reports the heat released per gram of whole samples, which means ash is also included. Ash is an inert mineral matter that adds mass without contributing anything to combustion. Thus, dividing the measured value by the combustible fraction $(1 - \text{ash content})$ removes this ballast. Assuming that ash releases no heat, the ash-free calorific value climbs with pyrolysis temperature, from 21.2\,MJ\,kg$^{-1}$ for the digestate to 25.1, 27.9, and 28.2\,MJ\,kg$^{-1}$ at 400, 500, and 600\,$^{\circ}$C respectively. The organic fraction of the biochar becomes more energy-dense as temperature increases, consistent with the progressive loss of oxygen-rich volatiles and the enrichment in fixed carbon that accompanies carbonization \parencite{pecchi-2019, hung-2017}. 

The high ash content additionally determines where these biochars are applicable. Half the mass of the 600\,$^{\circ}$C and 500\,$^{\circ}$C biochar is mineral residue, which makes it a poor option for solid fuel. On the other hand, taking the soil-amendment route seems like a better option since the same fraction could carry the mineral nutrients concentrated from the digestate or the biochar and return them to the soil with the carbon. In this context, the ash works in the material's favor; its effect on adsorption is a separate question, addressed along with the BET surface-area results.

\vspace{0.3em}

\textcite{hung-2017} reports about 16.6\,MJ\,kg$^{-1}$ for a solid digestate, containing roughly 23\,wt\% ash. The 16.41\,MJ\,kg$^{-1}$ measured here at 22.6\,\% ash closely reproduces that number, which can lay the foundation for the upcoming comparison of the biochars. The biochars themselves sit near the bottom of the published range. Table~\ref{tab:hhv_digestate_comparison} spans 13--22\,MJ\,kg$^{-1}$ after pyrolysis: \textcite{opatokun-2015} report 13--16\,MJ\,kg$^{-1}$ for a food-waste digestate and \textcite{kujawska-2026} 19--22\,MJ\,kg$^{-1}$ for an agricultural one. The 14.1--14.2\,MJ\,kg$^{-1}$ measured here falls in the lower band rather than below the range, close to the value Opatokun reports for the comparably high-ash material. Ash-content plays an important role. At roughly 50\,\% for the 600\,$^{\circ}$C biochar, this places this material at the top of the 15--50\,wt\% band typical of digestates.

To prove the points a bit further, using the equation~\ref{eq:edf}, The energy densification can be achieved by dividing the biochar calorific value by that of the feedstock. Every biochar has an energy densification factor of 0.86, and that means one thing. The pyrolysis lowered the energy density of the solid instead of raising it. Lignocellulosic feedstocks treated over the same 400--600\,$^{\circ}$C window normally reach 1.4--1.8 \parencite{bridgwater-2011, antal-2003}. This difference could actually be expected. For high-ash digestates, the factor should stay below 1.4 and can fall to 1 or lower. The reason is that ash dilution outweighs carbon enrichment \parencite{hung-2017, opatokun-2015}, and the present material reaches the lower bound of that prediction. Applying Equation~\ref{eq:energy-yield} with the mass
yields of Section~\ref{subsec:results-pyrolysis} gives energy yields of 39.5,
37.2 and 36.3\,\% at 400, 500 and 600\,$^{\circ}$C. For comparison, lignocellulosic feedstocks yield energy around 50--75\,\%.
However, if compared on an ash-free basis, the same biochars read differently. Their 25.1--28.2\,MJ\,kg$^{-1}$ lies inside the 25--33\,MJ\,kg$^{-1}$ reported for lignocellulosic
biochar (Table~\ref{tab:hhv_fuels}), and when the densification factor is recalculated again, this basis rises to 1.18--1.33. The ash-free biochar left by pyrolysis is comparable in energy density to wood-derived biochar, and the shortfall in the measured value comes from ash rather than from poor or incomplete carbonization. 
\newpage
\subsubsection{CHNS elemental analysis}
\setresponsible{Surya Salim}
\label{subsubsec:chns}

The results are evaluated on two bases. The as-analyzed basis reports the composition of the whole sample, mineral matter included. The ash-free basis divides each element by the combustible fraction $(1-\text{ash})$; therefore,describes the organic matter alone (Table~\ref{tab:chns-ashfree}). The Ash-content was determined for the three biochars in Section~\ref{subsubsec:calorific} but not for the activated and sulfur-loaded samples. This means that neither the oxygen, obtained by difference, nor any ash-free value can be calculated for the sulfur-loaded biochars. Because of this limitation, the pristine biochars are discussed on both bases, while the loaded biochars are evaluated through their as-analyzed contents and through the atomic ratios (Table~\ref{tab:chns-ratios}).


\begin{table}[ht]
\centering
\caption{CHNS elemental analysis of the French digestate biochars
and the corresponding activated, sulfur-loaded biochars, reported on the as-analyzed
basis.}
\label{tab:chns-full}
\begin{tabular}{llccccccc}
\toprule
\multirow{2}{*}{\thd{Sample}} & \multirow{2}{*}{\thd{$T$ ($^\circ$C)}}
& \multicolumn{4}{c}{\thd{Elemental content (\%)}}
& \multirow{2}{*}{\thd{Ash (\%)}}
& \multirow{2}{*}{\thd{O$_{\text{diff}}$ (\%)}}
& \multirow{2}{*}{\thd{Sum (\%)}} \\
\cmidrule(lr){3-6}
& & \thd{C} & \thd{H} & \thd{N} & \thd{S} & & & \\
\midrule
\multirow{3}{*}{\begin{tabular}[c]{@{}l@{}}Pristine\\biochar\end{tabular}}
& 400 & 37.98 & 1.534 & 2.744 & 1.486 & 43.5 & 12.76 & 100.0 \\
& 500 & 38.43 & 0.814 & 2.320 & 1.629 & 49.2 & 7.61  & 100.0 \\
& 600 & 37.99 & 0.521 & 1.626 & 1.584 & 50.0 & 8.28  & 100.0 \\
\midrule
\multirow{3}{*}{\begin{tabular}[c]{@{}l@{}}Sulfur-loaded\\activated biochar\end{tabular}}
& 400 & 31.66 & 0.342 & 1.148 & 6.804 & N/A & N/A & -- \\
& 500 & 33.61 & 0.344 & 1.066 & 5.290 & N/A & N/A & -- \\
& 600 & 33.47 & 0.348 & 0.943 & 5.217 & N/A & N/A & -- \\
\bottomrule
\end{tabular}

\vspace{2mm}

\end{table}

\begin{figure}[ht]
\centering
\begin{tikzpicture}
\begin{axis}[
    thesisbar, barlabels,
    title={Carbon content of the pristine biochars},
    bar width=20pt,
    symbolic x coords={400,500,600},
    xticklabels={{400\,$^{\circ}$C},{500\,$^{\circ}$C},{600\,$^{\circ}$C}},
    xlabel={Pyrolysis temperature},
    ylabel={Carbon content (\%)},
    ymax=90,
]
\addplot+[b primary] coordinates {(400,37.98) (500,38.43) (600,37.99)};
\addlegendentry{As-analyzed}
\addplot+[b emphasis] coordinates {(400,67.2) (500,75.7) (600,76.0)};
\addlegendentry{Ash-free}
\end{axis}
\end{tikzpicture}
\caption{Carbon content of the pristine biochars both as-analyzed and ash-free.}
\label{fig:chns-carbon}
\end{figure}


Figure~\ref{fig:chns-carbon} separates the carbon content on the two bases, as-analyzed and ash-free numbers. On the as-analyzed basis, the carbon content barely shows any difference in all three pyrolysis temperatures, sitting at 38.0, 38.4, and 38.0\,\% at 400, 500, and 600\,$^{\circ}$C. If plainly interpreted, that flatness would suggest that temperature did not have any influence to the carbon content. If we take the rising ash-content into account, which climbs from 43.5 to 50.0\,\% across the same range (Section~\ref{subsubsec:calorific}) and dilutes every organic element in equal measure, it shows a different picture. After dividing the carbon by the combustible fraction, which strips out that ballast, the ash-free carbon then rises from 67.2 to 75.7 to 76.0\,\%. The most noticeable gain is in the  400--500\,$^{\circ}$C step. The rise above 500\,$^{\circ}$C is very small. This means that the majority of carbonization happened below 500\,$^{\circ}$C. These ash-free numbers fall in the range reported for digestate biochars, whose carbon enriches more gently compared to lignocellulosic biochars because anaerobic digestion has already removed the fragile, easily devolatilized fractions before pyrolysis begins \cite{kujawska-2026}. 
Looking by the as-analyzed number, the loaded biochars visibly carry less carbon than the pristine biochars at 31-33\,\% against roughly 38\,\%. This is expected from the loading process, where the added sulfur, together with any iron it brings, adds mass that dilutes every element. 


\begin{figure}[ht]
\centering
\begin{tikzpicture}
\begin{axis}[
    thesissquare,
    title={H/C vs O/C : pristine biochar},
    width=0.78\textwidth,
    xlabel={Atomic O/C ratio}, ylabel={Atomic H/C ratio},
    xmin=0, xmax=0.45, ymin=0, ymax=0.66,
    xtick={0,0.1,0.2,0.3,0.4},
    ytick={0,0.1,0.2,0.3,0.4,0.5,0.6},
    legend columns=1,
    legend style={at={(0.03,0.97)}, anchor=north west, font=\small,
                  draw=tcGridGrey, fill=white},
]
\addplot[s primary] coordinates {(0.252,0.481) (0.149,0.252) (0.164,0.163)};
\addlegendentry{This work (biochar)}
\addplot[s reference] coordinates {(0.379,0.565) (0.138,0.086)};
\addlegendentry{Digestate biochar, Kujawska et al.}
% point labels placed off the trend lines
\node[anchor=south, yshift=3pt, font=\footnotesize, text=tcBlue] at (axis cs:0.252,0.481) {400\,$^{\circ}$C};
\node[anchor=east, xshift=-4pt, font=\footnotesize, text=tcBlue] at (axis cs:0.149,0.252) {500\,$^{\circ}$C};
\node[anchor=east, xshift=-4pt, font=\footnotesize, text=tcBlue] at (axis cs:0.164,0.163) {600\,$^{\circ}$C};
\node[anchor=south, yshift=3pt, font=\footnotesize, text=tcGrey] at (axis cs:0.379,0.565) {400\,$^{\circ}$C};
\node[anchor=north, yshift=-3pt, font=\footnotesize, text=tcGrey] at (axis cs:0.138,0.086) {800\,$^{\circ}$C};
\end{axis}
\end{tikzpicture}
\caption{Van Krevelen diagram for the French digestate biochars, with the digestate-biochar trajectory of \textcite{kujawska-2026} shown for reference.}
\label{fig:chns-vankrevelen}
\end{figure}

With the van Krevelen diagram (Figure~\ref{fig:chns-vankrevelen}), one can utilize the atomic H/C and O/C ratios to explain the changes in the degree of biochar carbonization. This method is widely applied as an indicator of the stability and long-term environmental properties of biochar \cite{kumar-2024}. It could be seen in figure~\ref{fig:chns-vankrevelen} that both ratios fall as the pyrolysis temperature rises. H/C drops from 0.481 at 400\,$^{\circ}$C to 0.252 at 500\,$^{\circ}$C and 0.163 at 600\,$^{\circ}$C, while O/C moves from 0.252 to 0.149. The downward-left trend is the signature of aromatization, where the biochar sheds hydrogen- and oxygen-bearing groups and its carbon condenses into fused aromatic rings. The steepest move is between 400-500\,$^{\circ}$C, which corresponds with the increase in carbon in figure~\ref{fig:chns-carbon}. The reference line of \textcite{kujawska-2026} for a manure digestate biochar also runs along the same diagonal. This biochar carbonizes along the near-linear path typical of a digestate that anaerobic digestion has already partially stabilized.
Since the oxygen is obtained by difference and thus contains every error in the other channels, the small rightward step in O/C between 500 and 600\,$^{\circ}$C (0.149 to 0.164) falls within that case. 

\begin{figure}[ht]
\centering
\begin{tikzpicture}
\begin{axis}[
    thesisbar, barlabels,
    title={H/C Ratio of the Biochars},
    bar width=20pt,
    symbolic x coords={400,500,600},
    xticklabels={{400\,$^{\circ}$C},{500\,$^{\circ}$C},{600\,$^{\circ}$C}},
    xlabel={Parent pyrolysis temperature},
    ylabel={Atomic H/C ratio},
    ymax=0.56,
    every node near coord/.append style={/pgf/number format/precision=3},
]
\addplot+[b primary] coordinates {(400,0.481) (500,0.252) (600,0.163)};
\addlegendentry{Pristine Biochar}
\addplot+[b emphasis] coordinates {(400,0.129) (500,0.122) (600,0.124)};
\addlegendentry{activated, sulfur-loaded biochar}
\end{axis}
\end{tikzpicture}
\caption{Atomic H/C ratio of the pristine and the activated, sulfur-loaded biochars.}
\label{fig:chns-hc-loaded}
\end{figure}




Now that in figure~\ref{fig:chns-hc-loaded} both biochars are set side by side, it can be observed, that once the biochars have been steam-activated and sulfur-loaded, that trend is gone. All three sit at 0.129, 0.122, and 0.124, a common floor near 0.12-0.13 across all temperatures. It is safe to assume that the hydrogen is stripped in the 700\,$^{\circ}$C activation step, not during sulfur loading. The size of the change tracks how much hydrogen was there to lose. The 400\,$^{\circ}$C biochar falls the furthest, from 0.481 to 0.129, while the 600\,$^{\circ}$C parent only moves from 0.163 to 0.124. The pyrolysis had already carbonized it close to the same level. All three loaded and activated biochars end up below even the 600\,$^{\circ}$C pristine Biochar number. This means that activation carries aromatization past what pyrolysis alone reaches at 600\,$^{\circ}$C. Another thing to take from this is that the pyrolysis temperature, which matters a lot for the pristine Biochar, has little influence on the aromaticity of the final activated biochar once the activation step has leveled it out.
\clearpage
\begin{figure}[ht]
\centering
\begin{tikzpicture}
\begin{axis}[
    thesisbar, barlabels,
    title={N/C Ratio of the Biochars},
    bar width=20pt,
    symbolic x coords={400,500,600},
    xticklabels={{400\,$^{\circ}$C},{500\,$^{\circ}$C},{600\,$^{\circ}$C}},
    xlabel={Parent pyrolysis temperature},
    ylabel={Atomic N/C ratio},
    ymax=0.072,
    every node near coord/.append style={/pgf/number format/precision=3},
]
\addplot+[b primary] coordinates {(400,0.062) (500,0.052) (600,0.037)};
\addlegendentry{Pristine Biochar}
\addplot+[b emphasis] coordinates {(400,0.031) (500,0.027) (600,0.024)};
\addlegendentry{activated, sulfur-loaded biochar}
\end{axis}
\end{tikzpicture}
\caption{Atomic N/C ratio of the pristine and the activated, sulfur-loaded biochars.}
\label{fig:chns-nc}
\end{figure}


Slaughterhouse digestate is protein-rich, and the pristine Biochars hold a correspondingly large share of nitrogen: 2.74, 2.32, and 1.63\,\% at 400, 500, and 600\,$^{\circ}$C on an as-analyzed basis. From the N/C ratio in figure~\ref{fig:chns-nc}, nitrogen decreases as the pyrolysis temperature rises. Compared to carbon, nitrogen is observed to leave faster. The N/C ratio drops from 0.062 to 0.037 across the three temperatures, indicating that each step of carbonization removes proportionally more nitrogen than carbon.  Labile nitrogen in amide and amine groups volatilizes while the carbon condenses and remains behind. The same decline in nitrogen with temperature is reported across different papers for digestate biochars \cite{kujawska-2026, tomczyk-2020}
 
The loaded biochars sit lower again, with N/C values at 0.031, 0.027, and 0.024, close to half of the related biochar at every temperature. Due to the activation process, it drove off the residual nitrogen, much like the residual hydrogen. H/C collapsed onto a single floor, but N/C still maintains the downward trend of the parent biochar.

\clearpage
\begin{figure}[ht]
\centering
\begin{tikzpicture}
\begin{axis}[
    thesisbar, barlabels,
    title={Sulfur content of the Biochars},
    bar width=20pt,
    symbolic x coords={400,500,600},
    xticklabels={{400\,$^{\circ}$C},{500\,$^{\circ}$C},{600\,$^{\circ}$C}},
    xlabel={Parent pyrolysis temperature},
    ylabel={Sulfur content, as-analyzed (\%)},
    ymax=7.8,
    every node near coord/.append style={/pgf/number format/precision=2},
]
\addplot+[b primary] coordinates {(400,1.486) (500,1.629) (600,1.584)};
\addlegendentry{Pristine Biochar}
\addplot+[b emphasis] coordinates {(400,6.804) (500,5.290) (600,5.217)};
\addlegendentry{activated, sulfur-loaded biochar}
\end{axis}
\end{tikzpicture}
\caption{Sulfur content of the pristine and the activated, sulfur-loaded chars.}
\label{fig:chns-s}
\end{figure}

The data shown by the figure~\ref{fig:chns-s} can be compared against the gravimetric loading capacities in Section~\ref{subsec:results-sulfur}. The results from table~\ref{tab:sulfur_loading} were obtained independently from the mass each sample gained during loading. Treating the loaded sulfur as a deposited mass and taking the biochar's own sulfur as the baseline, the deposited-sulfur capacity per gram, $Q_\mathrm{S}$, follows from the measured loaded fraction $w_\mathrm{S,loaded}$ and the baseline $w_\mathrm{S,base}$:
\begin{equation}
Q_\mathrm{S} = \frac{w_\mathrm{S,loaded} - w_\mathrm{S,base}}{1 - w_\mathrm{S,loaded}}
\label{eq:s-capacity}
\end{equation}
With $w_\mathrm{S,base}\approx1.5\,\%$, Equation~\ref{eq:s-capacity} gives 57, 40 and 39\,mg\,g$^{-1}$ at 400, 500, and 600\,$^{\circ}$C, against the gravimetric 41, 36, and 44\,mg\,g$^{-1}$. These two measures tell a broadly consistent story but do not match point for point. The 400\,$^{\circ}$C parent is widely agreed to be a strong loader. They diverge most at 600\,$^{\circ}$C, where the mass gain ranks highest (44\,mg\,g$^{-1}$), but here, the sulfur analysis ranks among the lowest (5.2\,\%). The major reason for this difference is that the two methods count different things. The mass gain records the net change (or decrease) in weight, which does not take some variables into account, like moisture, oxidation, and handling alongside the deposited sulfur. CHNS, on the other hand, responds to sulfur alone. The 400\,$^{\circ}$C number makes this plain, since the CHNS-derived 57\,mg\,g$^{-1}$ exceeds the 41\,mg\,g$^{-1}$ actually weighed in; this can only mean that the biochar shed some mass even as sulfur built up. At 600\,$^{\circ}$C, the reverse also happens, the weight gain higher than the sulfur-content number.



\newpage
\subsubsection{BET analysis}
\setresponsible{Regen Wijaya}
\label{subsubsec:bet}
The BET analysis of pristine, as well as sulfur-loaded and activated French biochars, resulted in three parameters: the SSA, total pore volume, and mean pore diameter.  

\begin{figure}[htbp]
\centering
\begin{tikzpicture}
\begin{axis}[
    thesisbar, barlabels,
    title={Specific surface area of French biochars},
    bar width=22pt,
    symbolic x coords={400C, 500C, 600C},
    xticklabels={400\textdegree C, 500\textdegree C, 600\textdegree C},
    xlabel={Pyrolysis temperature},
    ylabel={Specific surface area in m$^2$/g},
    ymax=160,
    scaled y ticks=false,
    every node near coord/.append style={/pgf/number format/precision=2},
]
\addplot[b primary] coordinates {
    (400C, 0.780)
    (500C, 1.800)
    (600C, 31.970)
};
\addlegendentry{Pristine biochar}

\addplot[b emphasis] coordinates {
    (400C, 100.966)
    (500C, 146.089)
    (600C, 88.636)
};
\addlegendentry{Activated, sulfur loaded biochar}

\end{axis}
\end{tikzpicture}
\caption{SSA of French pristine and activated, sulfur loaded biochars across different pyrolysis temperatures.}
\label{fig:bet_surface_area}
\end{figure}

Figure~\ref{fig:bet_surface_area} shows specifically the SSA of both pristine and activated, sulfur loaded French biochars measured during the BET-analysis. Pristine biochars exhibited an increase in SSA with increasing pyrolysis temperature. The lowest SSA values for the pristine biochars were measured at 400 °C at around 0.78 m$^2$\,g$^{-1}$ and at 500°C at around 1.8 m$^2$\,g$^{-1}$. At 600°C, the SSA reached the highest number at 31.97 m$^2$\,g$^{-1}$. This may be attributed to the promotion of de-volatilization and the development of internal pores and porous structure with increasing temperature, thereby increasing the accessible surface area \parencite{tomczyk-2020,ghorbani-2022}.

\vspace{0.3em}

Compared to pristine biochar samples, activated, sulfur-loaded biochars exhibited a substantially larger specific surface area. The SSA increased significantly to 100.97 m$^2$\,g$^{-1}$ at 400 °C and to 146.09 m$^2$\,g$^{-1}$ at 500 °C, suggesting that more micropores may be unblocked and developed during steam activation. Furthermore, this suggests that the biochar produced at a pyrolysis temperature of 500°C provided a particularly favorable carbon structure for subsequent pore formation, especially micropores, during steam activation, since the highest SSA is observed at 500°C. However, the activated, sulfur-loaded biochar produced from the 400°C exhibited the highest sulfur content around 6.804\%, as shown in Figure~\ref{fig:chns-s}, indicating that the pore structure of activated, sulfur-loaded biochar at 400°C may be partially occupied and blocked by sulfur element, which can reduce the internal surface accessible to N$_2$ during BET analysis. 

\vspace{0.3em}

In contrast to the SSA of pristine biochars, the SSA observed at 600°C was the lowest among the SSA of other activated, sulfur-loaded biochars. This may indicate that the carbon structure of biochar pyrolyzed at 600°C was strongly and less reactive for further gasification and pore development during steam activation. 

\vspace{0.3em}

The SSA values of the biochars in this study are lower compared to the SSA values reported in the literature, as shown in Table~\ref{tab:biochar_ssa_selected}, even though the feedstock differs significantly. The SSA of pristine biochar produced at 400°C in this work is 20 times lower compared to the SSA of biochar produced from pig manure at the same pyrolysis temperature, which is around 15.6 m$^2$\,g$^{-1}$. This number is comparably achieved by the biochar produced in this work at 600°C. Furthermore, the SSA of activated biochar derived from pig manure at pyrolysis temperature of 700°C reported by Vamvuka and Raftogianni \cite{Vamvuka2021PigManure} is about 231.4 m$^2$\,g$^{-1}$, which is higher than the maximum SSA value of 146.09 m$^2$\,g$^{-1}$ at 500°C measured in this work. This may strongly be caused by the difference of feedstock composition, of which pig manure can contain residual lignocellulosic material from undigested feed, whereas slaughterhouse digestate is richer in animal-derived organic matter and minerals.

\vspace{0.3em}

In addition, the SSA of activated, sulfur loaded biochars in this work is significantly lower compared to the SSA of plant-derived activated biochars, including 1479~m$^2$\,g$^{-1}$ for olive-stone-derived material and 767~m$^2$\,g$^{-1}$ for vine-shoot-derived material. These big differences in SSA are attributed particularly to the feedstock characteristics. As discussed in Section~\ref{Properties of Biochar}, plant-derived feedstocks generally contain a larger carbonaceous and lignocellulosic fraction and can develop a more extensive porous structure, whereas animal- and digestate-derived biochars usually contain a comparatively high mineral and ash fraction, which may reduce the accessible specific surface area \cite{Vamvuka2021PigManure} \cite{Gargiulo2018SteamBiochar}.

%As illustrated in Figure~\ref{fig:bet_surface_area}, the specific surface area (SSA) of the pristine biochar samples increased with rising pyrolysis temperature. The pristine biochar sample produced at pyrolysis temperatures of 400 °C and 500 °C exhibited a SSA of 0.78 and 1.8 m$^2$\,g$^{-1}$, respectively, which align with the range reported in the literature for animal-based biochar. However, the SSA of the pristine biochar sample at 600°C increased significantly to 31.97 m$^2$\,g$^{-1}$, suggesting that 




%The specific surface area of pristine biochar remains very low at 400°C and 500°C (0.78 and 1.8 m$^2$\,g$^{-1}$, respectively), indicating a largely non-porous or poorly developed pore structure at these temperatures. This is consistent with the fact that at lower pyrolysis temperatures, the carbonization process is incomplete — volatile matter and tars remain trapped within or block the developing pore network, limiting the accessible surface area. At 600°C, however, BC shows a sharp increase to 31.97 m$^2$\,g$^{-1}$. This jump reflects more complete pyrolysis: at higher temperatures, greater release of volatile organic compounds occurs, along with more advanced decomposition of the lignocellulosic structure, which opens up internal pores that were previously blocked or undeveloped. This suggests 600°C represents a threshold temperature at which significant pore development begins to occur naturally, even without any activation treatment.









%the pristine biochar samples exhibited BET specific surface areas of only 0.78, 1.80, and 31.97~m$^2$,g$^{-1}$ at pyrolysis temperatures of 400, 500, and 600~$^\circ$C, respectively. Thus, increasing the pyrolysis temperature alone already promoted considerable pore development, particularly between 500 and 600~$^\circ$C. This behavior agrees with the mechanism discussed in Section~\ref{subsubsec:Effect of Temperature on Biochar}: increasing temperature promotes devolatilization, whereby volatile compounds are progressively removed from the carbon matrix and previously inaccessible voids and pores become exposed. As a result, a more developed internal pore network, including micropores and mesopores, can form. Since micropores provide a particularly large internal surface per unit mass, their formation and exposure can strongly increase the BET specific surface area \parencite{tomczyk-2020,ghorbani-2022}.



%The effect of the subsequent treatment was substantially greater than that of pyrolysis alone. After steam activation and sulfur loading, the BET specific surface areas increased to 100.97, 146.09, and 88.64~m$^2$,g$^{-1}$ for the samples initially pyrolyzed at 400, 500, and 600~$^\circ$C, respectively. The corresponding absolute increases relative to the pristine biochars were approximately 100, 144, and 57~m$^2$,g$^{-1}$. These large differences provide strong evidence that the treatment opened a substantial fraction of previously inaccessible internal surface. During steam activation, H$_2$O penetrates into the pore system and partially gasifies the carbon matrix. This process can remove carbon from pore entrances, open blocked micropores, create new micropores, and enlarge existing pores into mesopores. Macropores contribute comparatively little directly to BET surface area because of their large dimensions, but they can serve as transport channels through which the adsorbate reaches smaller mesopores and micropores located deeper within the particle. The large increase in BET surface area therefore indicates the development of a more accessible hierarchical pore structure rather than simply the formation of larger pores \parencite{Sajjadi2019}.



%The most important finding is that the highest BET surface area was not obtained from the biochar pyrolyzed at the highest temperature. Instead, the 500~$^\circ$C precursor reached the maximum value of 146.09~m$^2$,g$^{-1}$ after treatment. This suggests that the 500~$^\circ$C biochar provided the most favorable initial structure for subsequent pore development. At this temperature, sufficient devolatilization had already occurred to establish an initial pore network, while the carbon matrix apparently retained enough reactivity for steam to further gasify pore walls and open additional micropores. The 500~$^\circ$C precursor can therefore be interpreted as representing a balance between initial carbonization and subsequent activation reactivity under the conditions investigated.



%The 400~$^\circ$C result further supports this interpretation. Although this sample exhibited a relatively high mass loss during steam activation, its final specific surface area reached only 100.97~m$^2$,g$^{-1}$, considerably below the 146.09~m$^2$,g$^{-1}$ obtained from the 500~$^\circ$C precursor. Therefore, activation mass loss cannot be interpreted directly as pore formation. The biochar initially produced at 400~$^\circ$C was less carbonized and consequently still contained a larger proportion of thermally unstable and volatile material. When reheated to the activation temperature of 700~$^\circ$C, part of the measured mass loss was likely associated with further devolatilization rather than selective carbon gasification responsible for creating new internal surface. In other words, a larger amount of material was removed, but this removal was apparently less efficient in generating BET-accessible microporosity.



%A different behavior is observed for the 600~$^\circ$C precursor. Before activation, this sample already possessed by far the largest BET surface area of the pristine biochars, at 31.97~m$^2$,g$^{-1}$. Nevertheless, after treatment its surface area increased only to 88.64~m$^2$,g$^{-1}$ and remained below those of both the treated 400 and 500~$^\circ$C samples. Together with the lower activation mass loss measured for the 600~$^\circ$C biochar, this indicates that the highly carbonized precursor was less susceptible to additional steam gasification. A substantial portion of its pore structure may already have been developed during pyrolysis, leaving less reactive carbon available for further pore creation during activation. 



%In addition, pore enlargement offers a plausible explanation for why stronger initial pore development does not necessarily lead to the largest final BET surface area. Steam gasification does not only create new micropores; continued carbon removal can also widen existing micropores into mesopores and, under stronger conversion, mesopores into macropores. While this process may increase total pore volume and improve transport through the particle, it does not necessarily increase specific surface area. A large number of narrow micropores provides considerably more internal surface than the same pore volume distributed among larger mesopores or macropores. Consequently, the decrease from 146.09~m$^2$,g$^{-1}$ for the 500~$^\circ$C precursor to 88.64~m$^2$,g$^{-1}$ for the 600~$^\circ$C precursor may indicate that pore development shifted from the generation of new high-surface-area micropores toward pore widening, or that parts of the pore network became less accessible \parencite{Sajjadi2019}.



%The BET results therefore complement the activation mass-loss results and reveal information that mass loss alone cannot provide. High burn-off does not automatically correspond to high surface area; the decisive factor is how carbon removal modifies the pore network. Among the investigated conditions, the biochar initially pyrolyzed at 500~$^\circ$C produced the highest accessible surface area after treatment and therefore appears to provide the most favorable precursor structure for steam-induced pore development. This condition most likely promoted effective formation and opening of micropores while avoiding excessive pore widening into larger mesopores and macropores.



%It should nevertheless be emphasized that the BET measurement was performed on samples that had undergone both steam activation and subsequent sulfur loading. The measured surface areas therefore represent the final pore accessibility after the complete treatment sequence and cannot be attributed quantitatively to steam activation alone. Sulfur deposited during the loading process may occupy or partially block micropores and mesopores, thereby reducing the surface accessible to nitrogen during BET analysis. Consequently, the actual surface area directly after steam activation may have been higher than the final values reported here. Separate BET measurements before and after sulfur loading would be required to distinguish pore generation during activation from subsequent pore occupation by sulfur.



%Overall, the results demonstrate that the development of specific surface area is governed not simply by the highest pyrolysis temperature or the greatest activation mass loss, but by the interaction between the initial carbon structure and its reactivity during steam activation. Under the experimental conditions used in this work, the biochar pyrolyzed at 500~$^\circ$C showed the most favorable balance, producing the largest BET-accessible internal surface and therefore the greatest apparent potential for adsorption-related applications.

\begin{figure}[H]
\centering
\begin{tikzpicture}
\begin{axis}[
    thesisbar, barlabels,
    title={Total pore volume of French biochars},
    bar width=22pt,
    symbolic x coords={400C, 500C, 600C},
    xticklabels={400\textdegree C, 500\textdegree C, 600\textdegree C},
    xlabel={Pyrolysis temperature},
    ylabel={Total pore volume in cm$^3$/g},
    ymax=0.1,
    ytick={0,0.02,0.04,0.06,0.08,0.1},
    yticklabels={0,0.02,0.04,0.06,0.08,0.1},
    scaled y ticks=false,
    every node near coord/.append style={/pgf/number format/precision=3},
]
\addplot[b primary] coordinates {
    (400C, 0.003)
    (500C, 0.005)
    (600C, 0.028)
};
\addlegendentry{Pristine biochar}

\addplot[b emphasis] coordinates {
    (400C, 0.071)
    (500C, 0.095)
    (600C, 0.068)
};
\addlegendentry{Activated, sulfur loaded biochar}

\end{axis}
\end{tikzpicture}
\caption{Total pore volume of French pristine and activated, sulfur loaded biochars across different pyrolysis temperatures.}
\label{fig: total pore volume}
\end{figure}

Figure~\ref{fig: total pore volume} shows the total pore volume of both the pristine and the activated, sulfur-loaded French biochars, as measured by BET analysis. The pattern is identical to that of the SSA. The pore volume of the pristine samples also increased with rising pyrolysis temperature, with the highest total pore volume observed at 600 °C at approximately 0.028 cm$^3$\,g$^{-1}$, followed by 0.005 cm$^3$\,g$^{-1}$ at 500 °C and 0.003 cm$^3$\,g$^{-1}$ at 400 °C. As the temperature rose, further pore networks formed as a result of devolatilization \parencite{tomczyk-2020,ghorbani-2022}.
\vspace{0.3em}

After activation and sulfur loading, the pore volumes increased significantly to 0.071~cm$^3$\,g$^{-1}$ at 400°C, 0.095~cm$^3$\,g$^{-1}$ at 500°C and 0.068~cm$^3$\,g$^{-1}$ at 600°C. The biochar pyrolyzed at 500°C remains favorable to obtain higher pore volume through steam activation, while the biochar pyrolyzed at 600°C exhibited the lowest total pore volume.

\vspace{0.3em}

\begin{figure}[H]
\centering
\begin{tikzpicture}
\begin{axis}[
    thesisbar, barlabels,
    title={Mean pore diameter of French biochars},
    bar width=22pt,
    symbolic x coords={400C, 500C, 600C},
    xticklabels={400\textdegree C, 500\textdegree C, 600\textdegree C},
    xlabel={Pyrolysis temperature},
    ylabel={Mean pore diameter in nm},
    ymax=20,
    scaled y ticks=false,
    every node near coord/.append style={/pgf/number format/precision=3},
]
\addplot[b primary] coordinates {
    (400C, 17.814)
    (500C, 10.933)
    (600C, 3.508)
};
\addlegendentry{Pristine biochar}

\addplot[b emphasis] coordinates {
    (400C, 2.804)
    (500C, 2.603)
    (600C, 3.076)
};
\addlegendentry{Activated, sulfur loaded biochar}

\end{axis}
\end{tikzpicture}
\caption{Mean pore diameter of French pristine and activated, sulfur loaded biochars across different pyrolysis temperatures.}
\label{fig: mean pore}
\end{figure}

Figure~\ref{fig: mean pore} shows the mean pore diameter of both the pristine and the activated, sulfur-loaded French biochars, as calculated using Equation~\ref{Equation mean pore}. This trend differs from that of the SSA and the total pore volume since the mean pore diameter is derived from the ratio \(4V/A\). The mean pore diameter of pristine biochars decreased from 17.814 nm to 3.508 nm as the pyrolysis temperature, specific surface area, and total pore volume increased. This could indicate that the pore structure is increasingly dominated by smaller accessible pores, which may also include micropores. 

\vspace{0.3em}

After steam activation and sulfur loading, the mean pore diameters continued to decrease to 2.804 nm at 400°C, 2.603 nm at 500°C, and 3.076 nm at 600°C. The sample at 500 °C exhibited the smallest mean pore diameter as well as the largest specific surface area and pore volume, which could signify the most highly developed fine-pore structure among the samples examined. 

\newpage
\subsection{Comparison of H\texorpdfstring{$_2$}{2}O and CO\texorpdfstring{$_2$}{2} activation}
\setresponsible{Surya Salim}
\label{subsec:h2o-vs-co2}

The individual results in Sections~\ref{subsec:results-activation}, \ref{subsec:results-sulfur}, and \ref{sec:characterisation} each touched on the CO$_2$-activated biochar from the preceding project work. In this section, the threads between them are laid out and read as one comparison. Because of the same parameter used in the char activation process, these two different data sets can be placed side by side with reasonable confidence \cite{MartinezHuete2024Projektarbeit}. It is important to specify two boundaries right away. The comparison is not a controlled head-to-head on a single sample, but rather two slaughterhouse digestates from different sources. Additionally, the surface-area aspect of the comparison is based on published values and the sulfur-loading reaction rather than a direct measurement because no BET measurement was performed on the CO$_2$-activated biochar. A stark difference could be seen in the activation mass loss. The steam-activated char lost more mass than the CO$_2$-activated char at every pyrolysis temperature, and the gap was widest at 600\,$^{\circ}$C, where steam removed about 10--11\,\% compared to roughly 2\,\% for CO$_2$ (Figure~\ref{fig:co2-vs-steam}; Table~\ref{tab:co2-activation-projektarbeit}). The reaction formula covered in Section~\ref{subsubsec:physical-activation} leads to this. Under similar conditions, steam gasification progresses more quickly and achieves a higher carbon conversion because the H$_2$O molecule is smaller than CO$_2$ and diffuses more easily into the pore system \cite{Sajjadi2019, Zhang2021TyreActivation}. 

\vspace{0.3em}

The surface-area has the largest limitation for a comparison due to the missing CO$_2$ BET measurement. The steam-activated chars here reached 89--146\,m$^2$\,g$^{-1}$ (Figure~\ref{fig:bet_surface_area}), but there is no available number for the CO$_2$ char to compare against; therefore, it is difficult to illustrate the difference from this work alone. Published comparisons on a common feedstock could help in this case. \textcite{Zhang2021TyreActivation} activated a spent tyre char with both same agents and found steam gave the higher BET surface area at equal carbon conversion, about 667\,m$^2$\,g$^{-1}$ against 435\,m$^2$\,g$^{-1}$ for CO$_2$. Steam also developed both micropores and mesopores on the char, whereas CO$_2$ worked mainly within the micropore range. The same pattern also reported in the work from \textcite{Sajjadi2019}. On that basis, compared to CO$_2$ activation of the same material, the steam pathway used here would be predicted to create a wider pore network. Nevertheless, this is still an inference from the literature rather than a measurable outcome. 

\vspace{0.3em}

The sulfur-loading response is the one piece of the present dataset that speaks to this indirectly. The steam-activated chars took up roughly two to three times more sulfur than the CO$_2$-activated chars, 36--64\,mg\,g$^{-1}$ against 12--24\,mg\,g$^{-1}$ (Tables~\ref{tab:sulfur_loading} and~\ref{tab:sulfur_loading_co2}). A wider pore network with mesopores that boost the adsorption rate is a plausible reason. As already noted in Section~\ref{subsec:results-sulfur}, uptake on these high-ash chars also depends on the mineral content and surface chemistry rather than on surface area alone, so the sulfur data reinforce the surface-area comparison instead of standing in for it \cite{Huang2022DigestateH2S}.
 
\vspace{0.3em}

Considered altogether, the comparison points one way without exaggerating it. For this kind of digestate and for the sulfur-sorbent target of the work, steam activation takes the upper hand in effectiveness. It removed more mass, and the resulting char absorbed significantly more sulfur than its CO$_2$-activated counterpart. When a tight, strongly microporous structure is desired instead of the wider pore-size distribution that steam typically produces, CO$_2$ activation maintains its own advantage \cite{Sajjadi2019}. There are some drawbacks that also need to be considered in this conclusion, as stated before : the two feedstocks were not identical, the CO$_2$ char was never measured by BET, and the surface areas reported here are for the combined activation-and-loading sequence rather than for activation alone. 


\subsection{Application of the Biochars}
\setresponsible{Regen Wijaya}
\label{subsec:utilization}
Based on the results of the characterization of the physicochemical properties of biochar conducted in this study, the potential applications of the biochar produced are limited to use in adsorption, carbon storage, and soil amendment. The BET results show that the activated, sulfur-loaded biochar produced at a pyrolysis temperature of 500°C has the best-developed pore structure among all biochars produced in this work, with the highest SSA and total pore volume, which makes it the most favorable precursor for adsorption-related applications. The measured H$_2$S uptake, however, did not follow the surface area. In the 0.5\,mL\,g$^{-1}$ series the 500\,°C char took up the least sulfur of the three (36.4\,mg\,g$^{-1}$, against 40.9 at 400\,°C and 43.9 at 600\,°C, Table~\ref{tab:sulfur_loading}). In addition, the CHNS sulfur content is highest for the 400\,°C char at 6.80\,\%. As discussed in Section~\ref{subsubsec:bet}, the deposited sulfur may block part of the pore system and reduce the surface accessible to N$_2$, so the 500\,°C biochar is named here for its pore structure rather than for a demonstrated advantage in sulfur loading.


\vspace{0.3em}

Pristine biochar produced at 400°C may contribute to soil fertility, since it contains the highest nitrogen at around 2.744\%, whereas the pristine biochar produced at 600°C might be promising for long-term carbon storage in soil, since it has the lowest H/C ratio, which indicates a stable and highly carbon-rich structure attributed to an increased proportion of aromatic structure, which is highly resistant to microbial degradation in the soil \cite{Li2023Biochar}. 

\vspace{0.3em}

The application of biochars as solid fuel may unfortunately less promising, as the highest measured gross calorific values of the biochars were only approximately 14.1--14.2~MJ/kg, mainly because of the high ash contents of 43.5--50.0\%, compared to commercial charcoal with heating values over 27 MJ/kg \parencite{Mencarelli2025Charcoal}, whereas commercial wood pellets exhibit generally higher calorific values of approximately 18.5--19.3 MJ/kg \parencite{Mack2025Pellets}.

\vspace{0.3em}

%Overall, the results indicate that no single pyrolysis temperature is optimal for every potential application. The 500,$^\circ$C biochar appears to provide the most favorable precursor for subsequent activation and adsorption-related use, whereas the 600,$^\circ$C biochar may be more suitable when a highly carbonized and stable material is required. Among the investigated utilization pathways, H$_2$S removal represents the most directly supported application, while the BET and elemental results additionally indicate potential for broader environmental adsorption and carbon-storage applications.

%Overall, the present results suggest that higher-value applications based on adsorption, gas purification, or carbon storage may be more advantageous for this slaughterhouse-digestate-derived biochar.


\subsection{Overall assessment and viability}
\setresponsible{Regen Wijaya}
\label{subsec:assessment}
With all of the results, analyses, and discussions, this section weighs whether the thesis examined here, turning slaughterhouse digestate into a steam-activated biochar, is worth pursuing beyond this work. The discussion here will remain qualitative, which means it rests on the material properties measured rather than on cost figures. For that, it would then need pilot-scale data that the present experiments cannot provide several results argue in favor. 

\vspace{0.3em}

The feedstock is a residue from a slaughterhouse that normally carries a disposal cost, so using it is a form of waste valorization rather than a raw-material expense. The steam activation method worked: it raised the specific surface area from about 1--2\,m$^2$\,g$^{-1}$ in the pristine low-temperature chars to 146\,m$^2$\,g$^{-1}$ after treatment (Section~\ref{sec:characterisation}), and the resulting char took up sulfur at up to about 64\,mg\,g$^{-1}$, roughly two to three times the uptake of the CO$_2$ route tested previously (Section~\ref{subsec:h2o-vs-co2}). The sulfur results also point to a concrete use: the removal of H$_2$S from biogas, which is a low- to moderate-duty gas-cleaning task where a moderate surface area can already be useful.

\vspace{0.3em}

The limits are just as obvious. The surface areas achieved here are modest compared to commercial activated carbons, which typically reach several hundred to over a thousand m$^2$\,g$^{-1}$. There is less carbon available for pore development for an animal-derived digestate, mostly because of the high ash content. The process must also pay for two high-temperature stages, and the performance results come with the method limitations already noted.

\vspace{0.3em}
Taken together, the concept is worth pursuing, but as a specific proposition rather than a general one. Its main strength is not peak adsorption but the combination of a near-free waste feedstock and a performance that is adequate for a low- to moderate-duty task. A realistic role would be a low-cost desulfurization or soil-amendment material derived from slaughterhouse waste, not a replacement for high-grade activated carbon, as mentioned in the previous section (section~\ref{subsec:utilization}).

\subsection{Error Analysis}
\setresponsible{Surya Salim}
\label{subsec:error-analysis}
The limitations of this work, such as the limited feedstock, the condition of the activation setup, and the outsourced surface and elemental analyses, were already explained in section~\ref{subsec:limitations}. In this section, the measurement errors that affect how far the reported numbers can be pushed are gathered and summarized. Several key quantities are inferred rather than measured directly and therefore, carry the errors of their inputs. The activation mass loss (Section~\ref{subsec:methods-activation}) equals the steam-gasified carbon only if that reaction removes all the mass, whereas residual devolatilization of the sub-700\,$^{\circ}$C-activated chars, condensation of injected water, and the manual syringe injection all contribute. Because of this, rather than being a direct measurement, it is interpreted as an upper limit on carbon conversion. By the interpretation of CHNS results, the oxygen is taken by difference (Equation~\ref{eq:O_fraction}) and thus accumulates the error of every other channel. The ash-free values depend on the separately measured ash content, which was taken from the calorimetric residue rather than from a standard furnace ashing (\S~\ref{subsubsec:methods-ash}). The sulfur content read gravimetrically (Section~\ref{subsec:results-sulfur}) diverges from the CHNS value (Section~\ref{subsubsec:chns}), because the mass gain also accounts for moisture, oxidation, and handling losses, whereas CHNS responds only to sulfur. In the BET section, the mean pore diameter is computed as $4V/A$ (Equation~\ref{Equation mean pore}) from two BET inputs. The BET surface areas were, moreover, measured after sulfur loading rather than after activation alone. The deposited sulfur may block pores, and the "real" post-activation area can be higher than that reported in this work (Section~\ref{sec:characterisation}). The scatter in the data can only be quantified for the repeated measurements, such as the pyrolysis data (triplicate, standard deviation at most 0.58 percentage points, Section~\ref{subsec:results-pyrolysis}) and the calorific values and ash contents (triplicate, Figure~\ref{fig:calorific-value}). The CHNS values are a four-run mean whose individual results were not released by the laboratory (Section~\ref{subsubsec:methods-chns}) and the BET parameters are averages of only two measurements, therefore no scatter can be attached to either. 


